\section{Error Analysis}

\subsection{Human Alignment Verification (Human-as-Judge)}

\subsubsection{Experimental Setup}

This experiment examines the agreement
between the GPT-5.5 Judge's hallucination judgments
and those of a human annotator,
following the human alignment methodology
of \citet{zhouOmniI2CHolisticBenchmark2026}.
Evaluation metrics include Accuracy, F1, Cohen's $\kappa$
(chance-corrected agreement),
and Type Agreement
(the proportion of matching hallucination type labels
among samples that both human and judge classify as hallucinated).
$\kappa$ values are interpreted
following the Landis and Koch (1977) scale:
$<$0.20 slight, 0.21--0.40 fair, 0.41--0.60 moderate,
0.61--0.80 substantial, 0.81--1.00 almost perfect.

Forty samples are drawn from MathVista testmini CoT results
via stratified sampling
(10 per model, approximately 5 hallucinated and 5 non-hallucinated each,
seed 42),
prioritizing coverage of all three hallucination types
within the hallucinated group.
The annotation interface displays each image,
question, model response, and reference answer.
The annotator judges whether hallucination is present
and, if so, labels the type
(faithfulness, factuality, or logical).
A single-blind design
(model name and judge verdict hidden)
is used to avoid annotation bias.
\Cref{fig:annotate_interface} shows the annotation interface.

\begin{figure}[htbp]
  \centering
  \includegraphics[width=0.85\textwidth]{annotate.png}
  \caption{Human annotation interface used in the alignment experiment.
    Each screen shows the image, question, model response,
    and reference answer,
    with model name and judge verdict hidden to avoid bias.}
  \label{fig:annotate_interface}
\end{figure}

\subsubsection{Overall Results}

\Cref{tab:human_overall} reports the overall alignment metrics,
and \Cref{tab:human_confusion} shows the confusion matrix.

\begin{table}[htbp]
  \centering
  \caption{Human alignment: overall metrics.}
  \label{tab:human_overall}
  \begin{tabular}{lc}
    \toprule
    Metric & Value \\
    \midrule
    Accuracy       & 85.0\% (34/40) \\
    Precision      & 70.0\% (14/20) \\
    Recall         & 100.0\% (14/14) \\
    F1             & 0.8235 \\
    Cohen's $\kappa$ & 0.700 (Substantial) \\
    \bottomrule
  \end{tabular}
\end{table}

\begin{table}[htbp]
  \centering
  \caption{Human alignment: confusion matrix
    (MLLM Judge vs.\ human annotator).}
  \label{tab:human_confusion}
  \begin{tabular}{l cc}
    \toprule
     & Judge: Hallucinated & Judge: Not hallucinated \\
    \midrule
    Human: Hallucinated     & TP = 14 & FN = 0 \\
    Human: Not hallucinated & FP = 6  & TN = 20 \\
    \bottomrule
  \end{tabular}
\end{table}

\Cref{tab:human_per_model} reports per-model alignment.

\begin{table}[htbp]
  \centering
  \caption{Human alignment: per-model results.}
  \label{tab:human_per_model}
  \small
  \begin{tabular}{lccccc}
    \toprule
    Model & $N$ & Accuracy & F1 & $\kappa$ \\
    \midrule
    GPT-5.4-mini (CoT)     & 10 & 1.000 & 1.000 & 1.000 \\
    Qwen3.5-35B-A3B (CoT)  & 10 & 1.000 & 1.000 & 1.000 \\
    Qwen3-VL-235B (CoT)    & 10 & 0.800 & 0.750 & 0.600 \\
    Gemini 2.5 Flash (CoT) & 10 & 0.600 & 0.333 & 0.200 \\
    \bottomrule
  \end{tabular}
\end{table}

Among the 14 samples where both human and judge agree
on the presence of hallucination,
Type Agreement reaches 92.9\% (13/14).
\Cref{tab:human_type_matrix} shows
the hallucination type confusion matrix.

\begin{table}[htbp]
  \centering
  \caption{Hallucination type agreement matrix
    (14 concordant hallucinated samples).}
  \label{tab:human_type_matrix}
  \begin{tabular}{l ccc}
    \toprule
    Human $\backslash$ Judge & faithfulness & factuality & logical \\
    \midrule
    faithfulness & 7 & 0 & 0 \\
    factuality   & 0 & 4 & 0 \\
    logical      & 1 & 0 & 2 \\
    \bottomrule
  \end{tabular}
\end{table}

Overall alignment is strong ($\kappa = 0.70$),
and the MLLM Judge does not miss
any human-identified hallucination (Recall = 100\%),
confirming that its detection sensitivity fully covers
the human standard.
GPT-5.4-mini and Qwen3.5-35B-A3B align perfectly with the human
($\kappa = 1.0$),
as these two models produce clear and stable hallucination patterns.
Gemini 2.5 Flash has the worst alignment ($\kappa = 0.20$),
with 4 of the 6 over-judgments attributed to this model.

All six over-judgments share a single error pattern:
age-estimation questions
where the model produces a reasonable refusal
(e.g., ``cannot determine age from appearance'')
that the MLLM Judge misclassifies as logical hallucination.
The root cause lies in the coupling
of the ``informativeness'' and ``hallucination'' dimensions
in the MMHal-Bench scoring rubric:
score 3 is the only slot matching
``low information + no hallucination,''
but because the model's refusal contains some information
(explaining the reason for refusal),
the judge tends to assign scores of 1--2,
pushing the sample into the ``hallucinated'' category.
This suggests that the scoring rubric lacks
an intermediate slot for ``moderately informative and hallucination-free''
responses.

\subsection{Judge Model Consistency Check}

This experiment examines whether evaluation conclusions
depend on the choice of judge model,
by comparing GPT-5.5 and Claude Opus 4.7
on the same sample set.
One hundred samples are randomly drawn
from MathVista testmini (seed 42),
yielding 100 responses per model (400 total).
Both judges use identical system prompts
(MMHal-Bench 0--6 rubric, five few-shot examples,
hallucination type classification).

\Cref{tab:judge_consistency} reports the results.

\begin{table}[htbp]
  \centering
  \caption{Judge consistency: GPT-5.5 vs.\ Claude Opus 4.7.}
  \label{tab:judge_consistency}
  \small
  \begin{tabular}{lc ccc cc cc}
    \toprule
    Model & $N$ & GPT HR & Claude HR & $\Delta$HR
      & GPT Avg & Claude Avg & $\kappa$ & Spearman $\rho$ \\
    \midrule
    GPT-5.4-mini     & 100 & 45.0\% & 44.0\% & $-$1.0pp & 3.27 & 3.01 & 0.980 & 0.805 \\
    Gemini 2.5 Flash & 100 & 37.0\% & 27.0\% & $-$10.0pp & 4.30 & 4.40 & 0.591 & 0.658 \\
    Qwen3.5-35B-A3B  & 100 & 14.0\% & 18.0\% & +4.0pp & 5.01 & 4.61 & 0.777 & 0.722 \\
    Qwen3-VL-235B    & 100 & 15.0\% & 22.0\% & +7.0pp & 4.91 & 4.38 & 0.638 & 0.669 \\
    \midrule
    Overall          & 400 & --- & --- & --- & --- & --- & 0.776 & 0.740 \\
    \bottomrule
  \end{tabular}
\end{table}

The overall Pearson $r$ is 0.7623 ($p < 0.001$).
The two judges show substantial overall agreement
($\kappa = 0.776$, $r = 0.762$),
indicating that evaluation results are reasonably robust
across judge models.
GPT-5.4-mini is the most stable ($\kappa = 0.98$),
while Gemini 2.5 Flash shows the largest divergence ($\kappa = 0.59$).
Despite per-model differences,
the relative model ranking is preserved under both judges:
both Qwen models perform best,
GPT-5.4-mini is in the middle,
and Gemini 2.5 Flash is the weakest.

\subsection{Threshold Sensitivity Analysis}

This analysis examines how sensitive
the Hallucination Rate is to the threshold $T = 3$.
For each threshold $T$ from 0 to 6,
HR($T$) is computed as $N_{\text{score}<T} / N_{\text{total}}$.

\Cref{tab:threshold_sensitivity} reports the results
for three representative thresholds.

\begin{table}[htbp]
  \centering
  \caption{Threshold sensitivity (MathVista Direct).}
  \label{tab:threshold_sensitivity}
  \begin{tabular}{l ccc cc}
    \toprule
    Model & HR($T$=2) & HR($T$=3) & HR($T$=4)
      & $\Delta(3{\to}2)$ & $\Delta(4{\to}3)$ \\
    \midrule
    GPT-5.4-mini     & 18.1\% & 45.1\% & 45.3\% & $-$27.0pp & +0.2pp \\
    Gemini 2.5 Flash & 1.2\%  & 34.3\% & 34.6\% & $-$33.1pp & +0.3pp \\
    Qwen3.5-35B-A3B  & 0.8\%  & 13.3\% & 13.3\% & $-$12.5pp & +0.0pp \\
    Qwen3-VL-235B    & 2.1\%  & 16.3\% & 16.4\% & $-$14.2pp & +0.1pp \\
    \bottomrule
  \end{tabular}
\end{table}

\begin{figure}[htbp]
  \centering
  \includegraphics[width=0.75\textwidth]{fig3_threshold_sensitivity_line.pdf}
  \caption{Hallucination rate as a function of threshold $T$
    (MathVista Direct).
    The vertical dashed line marks the official threshold $T = 3$.}
  \label{fig:threshold_sensitivity}
\end{figure}

HR($T$=3) and HR($T$=4) are nearly identical
(difference $<$0.3pp),
indicating that very few samples receive score 3
and the judge's scores follow a clear bimodal distribution.
The largest jump occurs between $T$=2 and $T$=3
(12.5--33.1pp),
reflecting a concentration of samples at score 2.
The model ranking is fully preserved
across all three thresholds,
confirming that the core conclusions
are insensitive to the threshold choice.

\subsection{Length Bias Analysis}

This analysis examines whether the MLLM Judge
exhibits systematic length bias.
Spearman correlation coefficients
between response length (word count) and judge score
are computed for all MathVista samples
in both Direct and CoT modes.

\Cref{tab:length_bias} reports the results.

\begin{table}[htbp]
  \centering
  \caption{Spearman correlation between response length
    and judge score (MathVista).
    $^{**}$: $p < 0.01$; $^{***}$: $p < 0.001$.}
  \label{tab:length_bias}
  \small
  \begin{tabular}{l cc cc cc}
    \toprule
    Model & Direct $\rho$ & CoT $\rho$
      & Direct median & CoT median
      & Direct Q1$\to$Q4 & CoT Q1$\to$Q4 \\
    \midrule
    GPT-5.4-mini
      & +0.086$^{**}$ & +0.015
      & 1 & 43
      & 3.19$\to$3.46 & 4.17$\to$3.95 \\
    Gemini 2.5 Flash
      & $-$0.142$^{***}$ & $-$0.218$^{***}$
      & 242 & 228
      & 4.61$\to$3.80 & 5.30$\to$4.33 \\
    Qwen3.5-35B-A3B
      & +0.530$^{***}$ & +0.006
      & 42 & 122
      & 4.56$\to$5.33 & 5.32$\to$5.26 \\
    Qwen3-VL-235B
      & +0.394$^{***}$ & $-$0.241$^{***}$
      & 40 & 103
      & 4.38$\to$4.92 & 5.32$\to$4.32 \\
    \bottomrule
  \end{tabular}
\end{table}

The MLLM Judge does not exhibit
a uniform ``longer $\to$ higher score'' bias.
The direction of correlation depends
on the quality pattern of each model's responses:
positive when longer responses carry more valid reasoning
(Qwen models in Direct mode),
negative when longer responses introduce more errors
(Gemini, Qwen3-VL under CoT).
CoT generally eliminates or reverses the positive correlation,
indicating that the judge is more sensitive
to noise in longer reasoning chains.

\subsection{Cross-Model Hallucination Consistency}

This analysis examines whether hallucinations overlap across models
on the same questions.
For each sample,
the number of models that produce hallucinations (0--4) is counted.

\Cref{tab:cross_model_consistency} reports the distribution.

\begin{table}[htbp]
  \centering
  \caption{Cross-model hallucination overlap.
    Hard Ratio = proportion of samples
    where $\geq$3 out of 4 models hallucinate.}
  \label{tab:cross_model_consistency}
  \begin{tabular}{lc ccccc c}
    \toprule
    Dataset & $N$ & 0/4 & 1/4 & 2/4 & 3/4 & 4/4 & Hard Ratio \\
    \midrule
    MathVista Direct & 1000 & 387 & 321 & 175 & 49 & 68 & 11.7\% \\
    MathVista CoT    & 1000 & 487 & 232 & 132 & 77 & 72 & 14.9\% \\
    VQA-RAD          & 451  & 59  & 118 & 122 & 81 & 71 & 33.7\% \\
    \bottomrule
  \end{tabular}
\end{table}

\begin{figure}[htbp]
  \centering
  \includegraphics[width=0.75\textwidth]{fig5_radar_comparison.pdf}
  \caption{Radar chart comparing model performance
    across six dimensions (accuracy = $1 - \text{HR}$).
    Larger area indicates better overall capability.}
  \label{fig:radar_comparison}
\end{figure}

VQA-RAD is the most difficult dataset:
only 13.1\% of samples are answered correctly by all models,
and 33.7\% are answered incorrectly by three or more,
indicating substantial model-independent intrinsic difficulty.
MathVista Direct has an even distribution:
38.7\% fully correct and 6.8\% fully incorrect.
After CoT, the Hard Ratio increases
from 11.7\% to 14.9\%,
reflecting the polarization effect of CoT
that causes model errors to converge.
